\centerline{\bf \Huge Учимся красить Круги Эйлера}

\problem{\textbf{1}} В классе все увлечены математикой или информатикой. Сколько
человек в классе, если математикой занимаются 48 человек, информатикой 47, а математикой
и информатикой 28?

\problem{\textbf{2}} В классе 29 человек. 15 из них занимаются в музыкальном кружке, 21 — в математическом. Сколько человек посещают оба кружка, если известно, что только Вовочка(ученик класса) не ходит ни в один из двух кружков?

\problem{\textbf{3}} У Пети есть карточки, у которых каждая сторона — красная или синяя.
Карточек с синей стороной — 10, карточек с красной стороной — 12, карточек
с разными сторонами — 5. Сколько карточек, у которых стороны одного цвета? 

\problem{\textbf{4}} Во дворе стоят машины. Некоторые из них — москвичи, а остальные — жигули. Некоторые из машин красные, а остальные белые. Некоторые из машин новые, а остальные — старые. Известно, что красных москвичей — 3, новых москвичей — 4, а новых красных машин — 5. При этом старых белых москвичей — 2, новых белых жигулей — 1, а старых красных москвичей вообще ни одного. Сколько во дворе новых красных москвичей, если всего машин 21, а старых белых жигулей — 6?

\problem{\textbf{5}} Из 100 ребят, отправляющихся в детский оздоровительный лагерь, кататься на сноуборде умеют 30 ребят, на скейтборде — 28, на роликах — 42. На скейтборде и на сноуборде умеют кататься 8 ребят, на скейтборде и на роликах — 10, на сноуборде и на роликах — 5, а на всех трех — 3. Сколько ребят не умеют кататься ни на сноуборде, ни на скейтборде, ни на роликах? (В число умеющих кататься на сноуборде включены те, кто умеет кататься ещё на чём-либо, и так далее).

\problem{\textbf{6}} На пиратском корабле трудятся 67 морских разбойников. У 47 из них есть
ухо, у 35 — глаз, а у 23 счастливчиков есть и то, и другое. Новая инструкция
Профсоюза Работников Абордажного Крюка предписывает корабельному врачу
учитывать также наличие носа. Оказалось, что 20 пиратов имеют нос, 12 — и
нос, и ухо, 11 — и нос, и глаз, а 5 — все три органа. Сколько пиратов не имеют
ничего из перечисленного?

\problem{\textbf{7}} На прогулку пошли четвероклассники и пятиклассники. Все они были либо
босиком, либо в тапочках. Четвероклассников было 24, а босых учеников 16.
Обутых пятиклассников было столько же, сколько босых четвероклассников.
Сколько учеников гуляли? 


